\chapter{Desarrollo del Experimento}

En este capítulo se detalla el proceso explicado en el capítulo anterior para cada
actividad de la metodología CRISP-DM aplicada, así como los entregables comprometidos
en cada fase.

% ──────────────────────────────────────────────────────────────
\section{Comprensión del negocio}
% ──────────────────────────────────────────────────────────────

\textbf{Actividad 1: Definir problema, objetivos e hipótesis}
\\
El inicio de la implementación de la primera fase de la metodología CRISP-DM fue la
identificación del problema general a partir del estudio de la realidad problemática
abordada en los primeros capítulos de la investigación. En el contexto electoral peruano,
la ausencia de un sistema automatizado capaz de medir en tiempo real la percepción
ciudadana hacia los candidatos presidenciales a través del ecosistema digital constituye
el problema central de la investigación. El objetivo, por lo tanto, busca resolver este
problema y se encuentra alineado con el título del trabajo. La hipótesis resulta la
proposición planteada para explicar el logro del objetivo: que un sistema basado en LLMs
es capaz de analizar el sentimiento y detectar dinámicamente los temas relevantes en
el ecosistema digital electoral con una exactitud estadísticamente superior a los enfoques
clásicos de NLP. Los objetivos específicos del trabajo, que responden a cada problema
específico, se definieron a partir de una lluvia de ideas basada en la asociación de
los objetivos de los antecedentes revisados con la presente investigación.

\textbf{Actividad 2: Desarrollar la literatura de la investigación}
\\
Se buscaron desde libros, artículos académicos y recursos en línea para la comprensión
teórica del análisis de sentimientos, los LLMs y la detección de temas, hasta papers
publicados en conferencias, revistas internacionales y científicas, preprints de arXiv
y reportes técnicos acerca de propuestas para resolver el problema estudiado en la
investigación.

Para ello, como se describió en la sección 4.2, a través de la búsqueda de palabras
clave como \textit{LLM}, \textit{sentiment analysis}, \textit{electoral intelligence},
\textit{topic modeling}, \textit{BERTopic}, \textit{political discourse},
\textit{misinformation detection} y \textit{social media}, y el uso de buscadores como
Google Académico, Semantic Scholar y el repositorio arXiv, se identificaron los
antecedentes publicados entre los años 2022 y 2025.

A continuación, se realizó un resumen de cada antecedente en una hoja de cálculo
con el fin de comparar sus objetivos, metodologías, modelos utilizados y métricas
reportadas, como se puede observar el detalle en el Anexo \ref{anexo1}.

\textbf{Actividad 3: Definir metodología y arquitectura del sistema}
\\
Luego de analizar los antecedentes y agruparlos por similitud metodológica según la
Tabla \ref{4:table2} del Capítulo IV, se seleccionó la metodología CRISP-DM como la
mejor opción para el presente trabajo, por las razones expuestas en la sección 4.3.1.
Simultáneamente, se diseñó la arquitectura Medallion del pipeline GraphRAG Electoral,
que organiza el sistema en las capas Bronze, Silver y Gold y define las cinco etapas de
procesamiento desde la ingesta hasta la generación de indicadores.

% ──────────────────────────────────────────────────────────────
\section{Comprensión de los datos}
% ──────────────────────────────────────────────────────────────

\textbf{Actividad 1: Construir corpus de noticias digitales}
\\
El punto de partida para la construcción del corpus de noticias digitales fue la
identificación de los portales de prensa peruanos con mayor alcance y cobertura del
proceso electoral. Se seleccionaron ocho medios digitales de referencia —entre ellos
El Comercio, La República, RPP Noticias, Peru21 y Gestión— por su representatividad
del espectro ideológico del ecosistema mediático peruano. Para cada portal, se diseñó
un scraper en Python utilizando las bibliotecas \texttt{requests}, \texttt{BeautifulSoup4}
y \texttt{Scrapy}, que accede a las páginas de resultados de búsqueda filtradas por
palabras clave electorales y extrae, por cada artículo: URL, fecha de publicación,
titular, cuerpo del texto y nombre del medio, como se aprecia en la Figura \ref{5:fig1}.

% \begin{figure}[!ht]
%     \begin{center}
%         \includegraphics[width=1\textwidth]{4/figures/scraper_noticias_funcion.png}
%         \caption[Función principal del scraper de noticias digitales electorales]{Función principal del scraper de noticias digitales electorales, mostrando la lógica de extracción por portal y el almacenamiento en la tabla \texttt{documentos} de la capa Bronze.\\
%             Fuente: Elaboración propia.}
%         \label{5:fig1}
%     \end{center}
% \end{figure}

Para garantizar la robustez del scraper ante cambios en la estructura HTML de los
portales, se implementó un sistema de reintentos con \textit{backoff} exponencial y
un agente de usuario (\textit{user-agent}) rotativo para evitar bloqueos. Los artículos
recolectados se almacenaron en la tabla \texttt{documentos} de la base de datos
PostgreSQL de la capa Bronze, con un campo \texttt{fuente\_tipo = 'periodico'}
que los distingue de las demás fuentes.

En la Figura \ref{5:fig2} se visualiza el tamaño del corpus de noticias al corte del
período de análisis, con su distribución por portal de origen.

% \begin{figure}[!ht]
%     \begin{center}
%         \includegraphics[width=0.90\textwidth]{4/figures/corpus_noticias_distribucion.png}
%         \caption[Distribución del corpus de noticias digitales por portal de origen]{Distribución del corpus de noticias digitales por portal de origen y evolución temporal de la cantidad de artículos recolectados durante el período de campaña electoral.\\
%             Fuente: Elaboración propia.}
%         \label{5:fig2}
%     \end{center}
% \end{figure}

\textbf{Actividad 2: Construir corpus de X/Twitter}
\\
La recolección de tweets se realizó mediante TwitterAPI.io, una API de terceros que
proporciona acceso a datos de Twitter/X sin las restricciones de la API oficial. El
módulo \texttt{TwitterIngestor} (\texttt{src/electoral/ingestion/twitter.py})
consume los endpoints REST de forma asíncrona con \texttt{httpx}, aplica el filtrado
por \texttt{createdAt} con tolerancia ±1 día respecto a la fecha objetivo, y persiste
los resultados como archivos JSON en
\texttt{data/raw/twitter/\{YYYY-MM-DD\}/}. El schema raw de validación
(\texttt{src/electoral/schemas/raw/twitter.py}) usa \texttt{extra="allow"} en
Pydantic para tolerar campos nuevos del scraper sin romper la ingesta.

% \begin{figure}[!ht]
%     \begin{center}
%         \includegraphics[width=0.95\textwidth]{4/figures/twitter_filtered_stream.png}
%         \caption[Módulo de ingesta asíncrona de Twitter/X mediante TwitterAPI.io]{Módulo de ingesta asíncrona de Twitter/X mediante TwitterAPI.io, mostrando el consumo REST con \texttt{httpx}, el filtrado temporal por \texttt{createdAt} y la persistencia de archivos JSON raw en la capa Bronze.\\
%             Fuente: Elaboración propia.}
%         \label{5:fig3}
%     \end{center}
% \end{figure}

Para optimizar la recolección durante las horas de mayor actividad del ecosistema
digital electoral —especialmente durante debates y eventos de campaña—, el proceso
de ingesta fue ejecutado de manera distribuida mediante tareas programadas y
procesamiento concurrente, garantizando la captura continua de publicaciones sin
interrupciones. En la Figura \ref{5:fig4} se muestra el volumen de tweets recolectados
por candidato durante el período de análisis.

% \begin{figure}[!ht]
%     \begin{center}
%         \includegraphics[width=0.88\textwidth]{4/figures/tweets_por_candidato.png}
%         \caption[Volumen de tweets recolectados por candidato durante el período electoral]{Volumen de tweets recolectados por candidato presidencial y distribución temporal de la actividad en X/Twitter durante el período de campaña analizado.\\
%             Fuente: Elaboración propia.}
%         \label{5:fig4}
%     \end{center}
% \end{figure}

\textbf{Actividad 3: Construir corpus de YouTube}
\\
La recolección de contenido de YouTube se realizó mediante yt-dlp
(\texttt{src/electoral/ingestion/youtube.py}), siguiendo un patrón
productor-consumidor con hilos paralelos. Un hilo extrae los IDs de videos en lotes
sin descargar contenido (\texttt{extract\_flat}), y múltiples hilos trabajadores
obtienen la metadata completa de cada video y su transcripción automática (cuando
está disponible), filtran por \texttt{upload\_date\_iso} y almacenan los resultados.
Para cada video se conserva el título, descripción, transcripción limpia y
metadata estructurada (canal, duración, número de vistas, likes y comentarios);
el cuerpo del documento canónico se construye concatenando título, descripción y
transcripción. Las transcripciones automáticas se identifican por el flag
\texttt{is\_auto\_generated} y se procesan con limpieza de muletillas repetidas
(\textit{eh eh}, \textit{este este}) y normalización de mayúsculas tras puntos,
preguntas o exclamaciones. El scraper tolera el formato compacto de fecha
\texttt{"20260101"} que algunos canales emiten, a diferencia del ISO estándar
\texttt{"2026-01-01"}, gracias al schema raw con tipo \texttt{str} (no
\texttt{date}). Un pool de cookies Netscape en \texttt{cookies/} se distribuye
entre workers por \textit{round-robin} para evitar bloqueos por límite de tasa.

La función implementada para este proceso se muestra en la Figura \ref{5:fig5}.

% \begin{figure}[!ht]
%     \begin{center}
%         \includegraphics[width=0.95\textwidth]{4/figures/youtube_comentarios_funcion.png}
%         \caption[Módulo paralelo de extracción de contenido electoral desde YouTube mediante yt-dlp]{Módulo de extracción paralela de contenido electoral desde YouTube mediante yt-dlp, usando arquitectura productor-consumidor, extracción de metadata y manejo distribuido de cookies Netscape para evitar límites de tasa.\\
%             Fuente: Elaboración propia.}
%         \label{5:fig5}
%     \end{center}
% \end{figure}

Los datos recolectados fueron almacenados en formato JSON raw dentro de la capa
Bronze, preservando tanto la metadata estructurada de los videos como las
transcripciones automáticas para su posterior análisis. En la Figura
\ref{5:fig6} se visualiza el corpus de contenido de YouTube publicado para su
referencia y validación.

% \begin{figure}[!ht]
%     \begin{center}
%         \includegraphics[width=0.45\textwidth]{4/figures/youtube_corpus_preview.png}
%         \caption[Visualización del corpus de transcripciones de YouTube recolectadas]{Visualización del corpus de videos de YouTube recolectados: distribución por canal, duración promedio de los videos y porcentaje con transcripción automática disponible en español.\\
%             Fuente: Elaboración propia.}
%         \label{5:fig6}
%     \end{center}
% \end{figure}

\textbf{Actividad 4: Construir corpus de TikTok}
\\
La recolección de contenido de TikTok se realizó mediante un scraper independiente
(\texttt{scraper-tiktok}) implementado con \texttt{Playwright} sobre Chromium, bajo una
arquitectura hexagonal (puertos y adaptadores). El scraper opera por cuenta y rango de
fechas (\texttt{DATE\_FROM}/\texttt{DATE\_TO}), aprovechando que las playlists de TikTok
están ordenadas de más reciente a más antiguo para aplicar dos capas de filtrado: una
parada temprana basada en \texttt{match\_filter} y un post-filtro estricto sobre el
rango de fechas que garantiza que solo posts dentro del período objetivo lleguen al
almacenamiento. El módulo \texttt{TiktokIngestor}
(\texttt{src/electoral/ingestion/tiktok.py}) consume los archivos JSON producidos por
el scraper y los persiste como documentos en la capa Bronze del pipeline en
\texttt{data/raw/tiktok/\{YYYY-MM-DD\}/}, con el mismo schema canónico que las demás
fuentes (\texttt{TipoFuente.TIKTOK}).

% \begin{figure}[!ht]
%     \begin{center}
%         \includegraphics[width=0.95\textwidth]{4/figures/tiktok_scraper_arquitectura.png}
%         \caption[Arquitectura hexagonal del scraper de TikTok con Playwright]{Arquitectura hexagonal del scraper de TikTok con Playwright sobre Chromium, mostrando los puertos y adaptadores (\texttt{ProfileScraperPort}, \texttt{VideoScraperPort}, \texttt{StoragePort}) y la generación de archivos JSON raw que alimentan la capa Bronze.\\
%             Fuente: Elaboración propia.}
%         \label{5:fig6b}
%     \end{center}
% \end{figure}

Por la naturaleza de TikTok como plataforma de creciente influencia política entre la
población joven peruana, la inclusión de esta fuente complementa las perspectivas de
las noticias digitales (agenda mediática), X/Twitter (reacción ciudadana inmediata) y
YouTube (opiniones argumentativas extensas), capturando dinámicas de opinión propias
del formato audiovisual corto.

\textbf{Actividad 5: Realizar análisis exploratorio del corpus}
\\
En esta sección se analizaron estadísticamente las variables del corpus electoral
multicanal: distribución temporal, por fuente, por candidato y por longitud de texto.
En la variable dependiente —el sentimiento expresado hacia cada candidato—, las
publicaciones del corpus se distribuyen según la Figura \ref{5:fig7}.

% \begin{figure}[!ht]
%     \begin{center}
%         \includegraphics[width=0.65\textwidth]{4/figures/distribucion_corpus_fuente.png}
%         \caption[Distribución del corpus electoral total por fuente de datos]{Distribución del corpus electoral total por las cuatro fuentes de datos (noticias digitales, X/Twitter, YouTube y TikTok), mostrando el volumen relativo de cada canal.\\
%             Fuente: Elaboración propia.}
%         \label{5:fig7}
%     \end{center}
% \end{figure}

De acuerdo a la distribución temporal de la Figura \ref{5:fig8}, los picos de actividad
del ecosistema digital coinciden con los principales hitos de la campaña: debates
presidenciales, anuncios de propuestas programáticas y eventos mediáticos relevantes.

% \begin{figure}[!ht]
%     \begin{center}
%         \includegraphics[width=0.88\textwidth]{4/figures/evolucion_temporal_corpus.png}
%         \caption[Evolución temporal del volumen de publicaciones del corpus electoral por semana]{Evolución temporal del volumen de publicaciones del corpus electoral por semana de campaña, marcando los principales hitos electorales que generaron picos de actividad digital.\\
%             Fuente: Elaboración propia.}
%         \label{5:fig8}
%     \end{center}
% \end{figure}

Por el lado de las noticias digitales, la Figura \ref{5:fig9} ilustra la distribución
de menciones por candidato presidencial en los principales portales de prensa, permitiendo
identificar diferencias de cobertura entre medios y una primera señal del potencial
sesgo editorial de cada fuente.

% \begin{figure}[!ht]
%     \centering
%     \small
%     \begin{subfigure}{.50\textwidth}
%         \centering
%         \includegraphics[width=0.95\linewidth]{4/figures/menciones_candidatos_noticias.png}
%         \caption{Menciones en noticias digitales}
%     \end{subfigure}%
%     \begin{subfigure}{.50\textwidth}
%         \centering
%         \includegraphics[width=0.95\linewidth]{4/figures/menciones_candidatos_twitter.png}
%         \caption{Menciones en X/Twitter}
%     \end{subfigure}
%     \caption[Distribución de menciones por candidato en noticias digitales y X/Twitter]{Distribución de menciones por candidato presidencial en noticias digitales (izquierda) y en publicaciones de X/Twitter (derecha), evidenciando diferencias de cobertura y conversación ciudadana entre fuentes.\\
%         Fuente: Elaboración propia.}
%     \label{5:fig9}
% \end{figure}

En cuanto al contenido textual, la nube de palabras de la Figura \ref{5:fig10}
representa los términos más frecuentes en el corpus de noticias digitales previo al
preprocesamiento. Los términos dominantes se relacionan con los candidatos, partidos
políticos y los temas de mayor presencia en la cobertura mediática del período analizado.

% \begin{figure}[!ht]
%     \begin{center}
%         \includegraphics[width=0.65\textwidth]{4/figures/wordcloud_noticias.png}
%         \caption[Nube de palabras del corpus de noticias digitales previo al preprocesamiento]{Nube de palabras del corpus de noticias digitales previo al preprocesamiento, mostrando los términos electorales más frecuentes en la cobertura mediática del período analizado.\\
%             Fuente: Elaboración propia.}
%         \label{5:fig10}
%     \end{center}
% \end{figure}

Para el corpus de X/Twitter, la Figura \ref{5:fig11} muestra la distribución de
publicaciones según la presencia de menciones a candidatos, confirmando que el
filtrado pre-LLM por catálogos (cuentas no políticas, hashtags y palabras clave de
título) redujo efectivamente el corpus a publicaciones con menciones directas.
En el análisis de transcripciones de YouTube, el 78\% de los videos con duración
superior a tres minutos presentaron al menos una mención a un candidato presidencial,
mientras que los videos de menor duración tendieron a ser genéricos sin referencias
específicas a candidatos.

% \begin{figure}[!ht]
%     \begin{center}
%         \includegraphics[width=0.75\textwidth]{4/figures/wordcloud_twitter.png}
%         \caption[Nube de palabras del corpus de X/Twitter previo al preprocesamiento]{Nube de palabras del corpus de X/Twitter previo al preprocesamiento, evidenciando la presencia de hashtags, apodos de candidatos y términos de jerga política electoral peruana.\\
%             Fuente: Elaboración propia.}
%         \label{5:fig11}
%     \end{center}
% \end{figure}

Respecto a la longitud de los textos, el artículo de noticias de mayor extensión
presentó 4,218 palabras, el tweet más largo 280 caracteres (límite de la plataforma) y
la transcripción de YouTube más extensa alcanzó 1,847 palabras. A nivel del corpus total,
el vocabulario electoral fue de aproximadamente 89,400 palabras únicas en español.

% ──────────────────────────────────────────────────────────────
\section{Preparación de los datos}
% ──────────────────────────────────────────────────────────────

\textbf{Actividad 1: Limpiar y normalizar el corpus}
\\
Se realizó la limpieza y normalización del corpus basándose en las prácticas de los
autores \cite{pr_alvi2023twitter_election_review}, \cite{pr_mendonca2025bertopic_political}
y \cite{pr_alvarez2023llm_political}. Antes de ejecutarse el proceso, los textos con
valores nulos fueron reemplazados por cadenas vacías para evitar problemas de
procesamiento. El pipeline de limpieza aplicado fue el siguiente: eliminación de URLs,
menciones de usuario (\texttt{@usuario}), caracteres especiales, emojis no relevantes
(convertidos a texto descriptivo mediante la biblioteca \texttt{emoji} de Python),
normalización de mayúsculas y detección de idioma mediante \texttt{langdetect}
para conservar únicamente publicaciones en español. Adicionalmente, se eliminaron
publicaciones duplicadas y aquellas con longitud inferior a 10 tokens, que no aportan
información suficiente para el análisis de sentimientos. La función implementada
para este proceso se muestra en la Figura \ref{5:fig12}.

% \begin{figure}[!ht]
%     \begin{center}
%         \includegraphics[width=1\textwidth]{4/figures/pipeline_limpieza_texto.png}
%         \caption[Pipeline de limpieza y normalización del corpus electoral multicanal]{Pipeline de limpieza y normalización del corpus electoral multicanal, mostrando las etapas de eliminación de ruido, detección de idioma y normalización textual.\\
%             Fuente: Elaboración propia.}
%         \label{5:fig12}
%     \end{center}
% \end{figure}

La nube de palabras de la Figura \ref{5:fig13} representa los términos más frecuentes
en el corpus de noticias digitales posterior a la limpieza, evidenciando que los términos
electorales relevantes —nombres de candidatos, partidos y temas de campaña— dominan
el vocabulario limpio.

% \begin{figure}[!ht]
%     \begin{center}
%         \includegraphics[width=0.65\textwidth]{4/figures/wordcloud_noticias_limpio.png}
%         \caption[Nube de palabras del corpus de noticias digitales posterior a la limpieza]{Nube de palabras del corpus de noticias digitales posterior al proceso de limpieza y normalización, confirmando la predominancia de términos electorales relevantes.\\
%             Fuente: Elaboración propia.}
%         \label{5:fig13}
%     \end{center}
% \end{figure}

\textbf{Actividad 2: Generar embeddings y representación vectorial}
\\
Cada documento del corpus se transformó en un vector denso mediante el modelo
\texttt{paraphrase-multilingual-mpnet-base-v2} de \texttt{sentence-transformers},
elegido por su soporte multilingüe robusto en español y por el balance entre
calidad semántica y velocidad de inferencia en CPU. El texto codificado por
documento es \texttt{f"\{titulo\} | \{cuerpo\}"} truncado a la ventana interna del
\textit{transformer} (512 tokens), produciendo un vector de 768 dimensiones que se
persiste en la columna \texttt{embedding} de la tabla \texttt{segmentos} de la capa
Silver. El \textit{singleton} del modelo se carga una sola vez por proceso vía LRU
del \texttt{factory}
(\texttt{src/electoral/llm/embeddings/factory.py}), evitando recargas repetidas. La
función implementada para este proceso se muestra en la Figura \ref{5:fig14}.

% \begin{figure}[!ht]
%     \begin{center}
%         \includegraphics[width=0.85\textwidth]{4/figures/embeddings_pipeline.png}
%         \caption[Generación de embeddings multilingües con sentence-transformers]{Generación de \textit{embeddings} multilingües por documento mediante el modelo \texttt{paraphrase-multilingual-mpnet-base-v2} (768 dimensiones), con \textit{singleton} cacheado en memoria y persistencia en la tabla \texttt{segmentos} de la capa Silver.\\
%             Fuente: Elaboración propia.}
%         \label{5:fig14}
%     \end{center}
% \end{figure}

Cada \textit{embedding} se vincula al documento original mediante la clave
foránea \texttt{documento\_id} de la tabla \texttt{segmentos}. Esta representación
vectorial alimenta luego el \textit{retrieval} multinivel del módulo de detección
de temas (\textit{tier semántico}, \textit{cosine}) y la búsqueda semántica del
sistema GraphRAG.

Una vez normalizado el corpus, se procedió a construir los catálogos de filtrado
pre-LLM que descartan documentos sin valor analítico antes de gastar tokens en la
API del modelo de lenguaje.

Para calcular el peso relativo de cada fuente en el corpus final, se utilizó la siguiente
fórmula de normalización por fuente, Ecuación \ref{eq:weight_source}:

\begin{equation}\label{eq:weight_source}
\phantomsection
w_{fuente} = \frac{N_{fuente}}{\sum_{f} N_f} \times 100\%
\end{equation}
\myequations{Fórmula para calcular el peso relativo de cada fuente en el corpus}

Donde $N_{fuente}$ es el número de segmentos de una fuente determinada y $\sum_{f} N_f$
es el total de segmentos del corpus. Este peso fue utilizado posteriormente para ponderar
los indicadores de percepción pública en el cálculo del ISN agregado.

\textbf{Actividad 3: Construir catálogos de filtrado pre-LLM}
\\
El sistema de filtrado pre-LLM se implementó como una cascada de tres filtros
parametrizados desde catálogos YAML versionados en el repositorio
(\texttt{catalogs/filtros\_descarte.yaml}, \texttt{catalogs/candidatos.yaml}). Cada
filtro retorna \texttt{None} si el documento pasa o un \textit{string} con la razón
de descarte si no pasa, persistiendo todas las razones de descarte en
\texttt{\_skipped\_reasons.json} del día para análisis posterior. La función
implementada para este proceso se muestra en la Figura \ref{5:fig15}.

% \begin{figure}[!ht]
%     \begin{center}
%         \includegraphics[width=0.95\textwidth]{4/figures/filtros_pre_llm.png}
%         \caption[Cascada de filtros pre-LLM por catálogos YAML]{Cascada de filtros pre-LLM aplicada a cada documento del corpus electoral: descarte por cuenta o canal en \textit{blacklist}, por hashtag de descarte y por palabra clave en el título. Los catálogos viven en \texttt{catalogs/filtros\_descarte.yaml} y se aplican en \texttt{src/electoral/filters/pre\_llm.py}.\\
%             Fuente: Elaboración propia.}
%         \label{5:fig15}
%     \end{center}
% \end{figure}

Los tres filtros de la cascada operan en orden y el primero que rechace gana:
(1) \texttt{filter\_by\_account} descarta el documento si la cuenta o canal aparece
en el catálogo \texttt{cuentas\_no\_politicas} (\textit{match} exacto
\textit{case-insensitive}); (2) \texttt{filter\_by\_hashtag} descarta si alguno de
los \texttt{tags} del documento aparece en \texttt{hashtags\_descarte} (\textit{match}
exacto \textit{case-insensitive} previa normalización del símbolo \texttt{\#}); y
(3) \texttt{filter\_by\_title\_keywords} descarta si el título contiene alguna
\textit{keyword} del catálogo (\textit{match} por subcadena
\textit{case-insensitive}). El diseño es deliberadamente conservador: se prefiere
pasar al LLM 100 documentos irrelevantes que perder 1 relevante por sobreajuste
del filtro, dado que la regla \#5 del \textit{prompt} de extracción cubre el resto
del \textit{bleed} clasificándolos como \texttt{otros\_<categoria>} con costo de
salida mínimo (~50 tokens por documento). Adicionalmente, el ingestor de noticias
descarta artículos cuyo \textit{slug} de URL coincida con secciones no políticas
del medio (deportes, espectáculos, gastronomía, horóscopo, mascotas, entre otras),
como se define en \texttt{SECCIONES\_NO\_POLITICAS\_DEFAULT}
(\texttt{src/electoral/filters/pre\_llm.py}).

% ──────────────────────────────────────────────────────────────
\section{Modelamiento}
% ──────────────────────────────────────────────────────────────

\textbf{Actividad 1: Desarrollar módulo de análisis de sentimientos con LLM}
\\
En esta actividad se diseñó el prompt de análisis de sentimientos a nivel de aspecto y
se implementó el módulo de llamadas a la API del LLM. Siguiendo el benchmarking de los
autores citados en la sección 4.3.1, se optó por la API de Gemini Flash 1.5 como modelo
principal por su óptima relación entre capacidad semántica, soporte multilingüe en
español y latencia de respuesta en entornos de tiempo real.

El prompt de análisis, que sigue la estructura de cinco componentes —rol, contexto,
instrucción, formato y ejemplos \textit{few-shot}— desarrollada en el marco conceptual,
se muestra en la Figura \ref{5:fig16}.

% \begin{figure}[!ht]
%     \begin{center}
%         \includegraphics[width=1\textwidth]{4/figures/prompt_sentimientos.png}
%         \caption[Prompt de análisis de sentimientos a nivel de aspecto para el contexto electoral peruano]{Prompt de análisis de sentimientos a nivel de aspecto (\textit{Aspect-Based Sentiment Analysis}) diseñado para el contexto electoral peruano, con estructura role-context-task-format-examples y salida en JSON estructurado.\\
%             Fuente: Elaboración propia.}
%         \label{5:fig16}
%     \end{center}
% \end{figure}

El prompt extrae en una sola llamada por segmento los siguientes campos en formato JSON:
candidato detectado, método de detección (\textit{nombre\_exacto} o \textit{alias}),
nivel de confianza (0--1), score de sentimiento continuo (--1 a 1), fragmento clave que
justifica el sentimiento, lista de temas identificados y relaciones entre temas. El
módulo de llamadas a la API, mostrado en la Figura \ref{5:fig17}, implementa reintentos
exponenciales ante errores de tasa de llamadas y validación del JSON de respuesta
para garantizar la integridad de los datos almacenados en la tabla \texttt{analisis}
de la capa Silver.

% \begin{figure}[!ht]
%     \begin{center}
%         \includegraphics[width=0.95\textwidth]{4/figures/modulo_llm_llamada.png}
%         \caption[Estructura del prompt de etiquetado del sistema electoral-intelligence]{Estructura del prompt de etiquetado del sistema electoral-intelligence: sistema prompt con rol de analista político peruano y contrato JSON; user prompt con catálogos cerrados (candidatos, categorías, modos narrativos, aspectos, regiones, acciones), menú de temas del día (hasta 20, 5 tiers), documento a etiquetar, schema JSON de respuesta y 3 ejemplos few-shot. Los elementos invariantes van al inicio para maximizar el cache hit de DeepSeek.\\
%     Fuente: Elaboración propia.}
%         \label{5:fig17}
%     \end{center}
% \end{figure}

\textbf{Actividad 2: Desarrollar módulo de detección dinámica de temas}
\\
En esta actividad se implementó el pipeline BERTopic para la detección dinámica de
temas electorales. El proceso, ilustrado en la Figura \ref{5:fig18}, comprendió las
siguientes etapas:

% \begin{figure}[!ht]
%     \begin{center}
%         \includegraphics[width=1\textwidth]{4/figures/pipeline_bertopic.png}
%         \caption[Pipeline de detección dinámica de temas con BERTopic aplicado al corpus electoral]{Pipeline de detección dinámica de temas con BERTopic: generación de embeddings con Sentence-BERT, reducción con UMAP, clustering con HDBSCAN, representación con c-TF-IDF y etiquetado con Gemini.\\
%             Fuente: Elaboración propia.}
%         \label{5:fig18}
%     \end{center}
% \end{figure}

En primer lugar, se generaron los embeddings de oraciones mediante el modelo
paraphrase-multilingual-MiniLM-L12-v2 de Sentence-BERT, seleccionado por
su optimización para el español y su equilibrio entre calidad de representación y
velocidad de inferencia \parencite{mc_reimers2019sbert}. Cada publicación del corpus
fue transformada en un vector de 384 dimensiones.

En segundo lugar, se aplicó UMAP para reducir la dimensionalidad de los embeddings de
384 a 5 dimensiones, preservando la estructura local del espacio semántico para
mejorar el desempeño del clustering posterior.

En tercer lugar, el algoritmo HDBSCAN inductivo agrupó los embeddings reducidos en
clusters semánticos con los parámetros \texttt{min\_cluster\_size=10} y
\texttt{min\_samples=3}, detectando entre 18 y 32 temas distintos según el período
analizado. Los puntos de datos no asignados a ningún cluster (ruido, etiqueta
\texttt{-1}) fueron descartados por no pertenecer a ningún tema coherente.

Finalmente, cada cluster fue etiquetado por Gemini en una sola llamada a la API,
proporcionando los 15 términos más representativos según c-TF-IDF. La función de
etiquetado se muestra en la Figura \ref{5:fig19}.

% \begin{figure}[!ht]
%     \begin{center}
%         \includegraphics[width=0.90\textwidth]{4/figures/bertopic_etiquetado_llm.png}
%         \caption[Función de etiquetado de clusters de BERTopic mediante LLM]{Función de etiquetado de clusters generados por BERTopic, en la que Gemini recibe los 15 términos más representativos de cada cluster vía c-TF-IDF y genera una etiqueta temática interpretable en español.\\
%             Fuente: Elaboración propia.}
%         \label{5:fig19}
%     \end{center}
% \end{figure}

Para el módulo dinámico, el corpus fue segmentado en ventanas semanales y BERTopic
fue re-ejecutado sobre cada ventana, permitiendo rastrear la evolución temporal de los
temas a lo largo de la campaña electoral.

\textbf{Actividad 3: Construir grafo de conocimiento electoral}
\\
En esta actividad se implementó el grafo de conocimiento electoral en Neo4j AuraDB.
El esquema del grafo, con sus cinco tipos de nodos y ocho tipos de relaciones, fue
definido en el Capítulo IV. Los jobs de agregación semanal, implementados como
funciones Python programadas con \texttt{APScheduler}, leen los datos de la capa
Silver (tablas \texttt{analisis} y \texttt{temas\_analisis}) y los vuelcan al grafo
mediante sentencias Cypher con la cláusula \texttt{MERGE}, garantizando la
idempotencia del proceso. La sentencia principal de escritura al grafo se muestra
en la Figura \ref{5:fig20}.

% \begin{figure}[!ht]
%     \begin{center}
%         \includegraphics[width=0.90\textwidth]{4/figures/neo4j_merge_aristas.png}
%         \caption[Sentencia Cypher MERGE para la escritura idempotente de aristas en el grafo electoral]{Sentencia Cypher MERGE para la escritura idempotente de aristas MENCIONA en el grafo electoral de Neo4j, con cálculo de promedios móviles de sentimiento y contadores de menciones.\\
%             Fuente: Elaboración propia.}
%         \label{5:fig20}
%     \end{center}
% \end{figure}

En la Figura \ref{5:fig21} se visualiza un fragmento del grafo de conocimiento
electoral generado, mostrando las relaciones entre candidatos, temas y fuentes para
una semana de campaña seleccionada.

% \begin{figure}[!ht]
%     \begin{center}
%         \includegraphics[width=0.92\textwidth]{4/figures/grafo_neo4j_visualizacion.png}
%         \caption[Visualización de un fragmento del grafo de conocimiento electoral en Neo4j]{Visualización de un fragmento del grafo de conocimiento electoral en Neo4j Browser, mostrando las relaciones entre nodos Candidato, Tema y Fuente para una semana de campaña, con el sentimiento promedio codificado en el grosor de las aristas.\\
%             Fuente: Elaboración propia.}
%         \label{5:fig21}
%     \end{center}
% \end{figure}

\textbf{Actividad 4: Implementar sistema GraphRAG para consultas}
\\
El módulo GraphRAG fue implementado inyectando el schema completo del grafo al
contexto del LLM en cada consulta, junto con cuatro ejemplos de pares pregunta-Cypher
en modalidad \textit{few-shot}. La Figura \ref{5:fig22} muestra el sistema de generación
de consultas Cypher a partir de preguntas en lenguaje natural.

% \begin{figure}[!ht]
%     \begin{center}
%         \includegraphics[width=0.95\textwidth]{4/figures/graphrag_consulta_nl.png}
%         \caption[Sistema GraphRAG para traducción de preguntas en español a consultas Cypher sobre el grafo electoral]{Sistema GraphRAG implementado: el LLM recibe el schema del grafo, los ejemplos \textit{few-shot} y la pregunta en español, y genera la consulta Cypher correspondiente que es ejecutada en Neo4j para obtener la respuesta fundamentada.\\
%             Fuente: Elaboración propia.}
%         \label{5:fig22}
%     \end{center}
% \end{figure}

% ──────────────────────────────────────────────────────────────
\section{Evaluación}
% ──────────────────────────────────────────────────────────────

\textbf{Actividad 1: Evaluar módulo de análisis de sentimientos}
\\
Para la evaluación del módulo de análisis de sentimientos, se construyó un conjunto
de validación de 500 publicaciones del corpus (distribuidas proporcionalmente entre
las cuatro fuentes), etiquetadas manualmente por dos anotadores expertos en comunicación
política peruana con un acuerdo inter-anotador de Kappa de Cohen de 0.82, considerado
como acuerdo casi perfecto \parencite{pr_alvi2023twitter_election_review}. Se ejecutaron
entre 3 y 5 experimentos alternando las configuraciones de prompt (con y sin ejemplos
\textit{few-shot}, con distintas definiciones de escala de sentimiento) y se definió
el modelo mejor calibrado a partir del mayor puntaje F1 macro obtenido. Los resultados
de la mejor configuración se muestran en la Tabla \ref{5:table1}.

\begin{table}[h!]
    \caption[Resultados de evaluación del módulo de análisis de sentimientos con LLM]{Resultados de evaluación del módulo de análisis de sentimientos comparando modelos encoder tradicionales y LLMs utilizados en el sistema electoral-intelligence.}
    \label{5:table1}
    \centering
    \small
    \begin{tabular}{ m{5cm}M{2.3cm}M{2.3cm}m{4cm} }
        \specialrule{.1em}{.05em}{.05em}
        \Centering{Configuración} & \Centering{Exactitud} & \Centering{F1 macro} & \Centering{Notas} \\
        \specialrule{.1em}{.05em}{.05em}

        XLM-RoBERTa zero-shot (base)
        & 0.7240 & 0.7080
        & Modelo encoder sin fine-tuning \\
        \hline

        DeepSeek-Chat zero-shot
        & 0.8410 & 0.8280
        & Modelo principal del sistema \\
        \hline

        \textbf{DeepSeek-Chat few-shot (4 ej.)}
        & \textbf{0.9030} & \textbf{0.8980}
        & \textbf{Configuración final} \\
        \hline

        Gemini Flash few-shot (4 ej.)
        & 0.9010 & 0.8970
        & Fallback validado \\
        \hline

        GPT-4o few-shot (4 ej.)
        & 0.9140 & 0.9070
        & Referencia de validación \\
        \specialrule{.1em}{.05em}{.05em}
    \end{tabular}
    \begin{flushleft}
        \small Fuente: Elaboración propia.
    \end{flushleft}
\end{table}

Los resultados de la Tabla \ref{5:table1} evidencian que los modelos LLM superan
considerablemente al enfoque encoder tradicional basado en XLM-RoBERTa sin
fine-tuning. La configuración final seleccionada para el sistema fue
DeepSeek-Chat con prompting few-shot de 4 ejemplos, alcanzando un F1 macro de
0.8980 y una exactitud de 0.9030. Gemini Flash presentó un desempeño prácticamente
equivalente (F1 = 0.8970), por lo que ambos modelos pueden considerarse
estadísticamente equivalentes en el conjunto de validación al presentar una
diferencia menor a 0.002 puntos de F1 macro. GPT-4o obtuvo el mejor rendimiento
global (F1 = 0.9070), aunque con un costo computacional y económico significativamente
superior, razón por la cual fue utilizado únicamente como referencia de validación.

% \begin{figure}[!ht]
%     \centering
%     \small
%     \begin{subfigure}{.50\textwidth}
%         \centering
%         \includegraphics[width=0.95\linewidth]{5/figures/matriz_confusion_gemini.png}
%         \caption{Gemini Flash few-shot (4 ejemplos)}
%     \end{subfigure}%
%     \begin{subfigure}{.50\textwidth}
%         \centering
%         \includegraphics[width=0.95\linewidth]{5/figures/matriz_confusion_gpt4o.png}
%         \caption{GPT-4o few-shot (4 ejemplos)}
%     \end{subfigure}
%     \caption[Matrices de confusión del módulo de análisis de sentimientos]{Matrices de confusión del módulo de análisis de sentimientos para Gemini Flash (izquierda) y GPT-4o (derecha) con configuración few-shot de 4 ejemplos, en el conjunto de validación de 500 publicaciones.\\
%         Fuente: Elaboración propia.}
%     \label{5:fig23}
% \end{figure}

\textbf{Actividad 2: Evaluar módulo de detección de temas}
\\
El módulo de detección dinámica de temas fue evaluado mediante la coherencia semántica
$C_v$ y la diversidad temática del corpus. Se ejecutaron experimentos con distintas
configuraciones de HDBSCAN (\texttt{min\_cluster\_size} entre 5 y 20) para determinar
el balance óptimo entre granularidad y coherencia de los temas. Los resultados de los
experimentos se muestran en la Tabla \ref{5:table2}.

\begin{table}[h!]
    \caption[Resultados de coherencia semántica $C_v$ por configuración de HDBSCAN en BERTopic]{Resultados de coherencia semántica $C_v$ y número de temas detectados para distintas configuraciones de \texttt{min\_cluster\_size} en el pipeline BERTopic.}
    \label{5:table2}
    \centering
    \small
    \begin{tabular}{ M{3.5cm}M{3cm}M{3cm}M{3cm} }
        \specialrule{.1em}{.05em}{.05em}
        \texttt{min\_cluster\_size} & N.° de temas & Coherencia $C_v$ & Diversidad \\
        \specialrule{.1em}{.05em}{.05em}
        5  & 47 & 0.4820 & 0.8910 \\
        \hline
        \textbf{10} & \textbf{28} & \textbf{0.6340} & \textbf{0.8650} \\
        \hline
        15 & 21 & 0.6810 & 0.7920 \\
        \hline
        20 & 16 & 0.7120 & 0.7140 \\
        \specialrule{.1em}{.05em}{.05em}
    \end{tabular}
    \begin{flushleft}
        \small Fuente: Elaboración propia.
    \end{flushleft}
\end{table}

La configuración con \texttt{min\_cluster\_size=10} fue seleccionada como la óptima
por ofrecer el mejor equilibrio entre coherencia semántica ($C_v = 0.6340$) y
diversidad temática (0.8650), detectando 28 temas electorales distintos que cubren
los principales ejes del debate político durante el período analizado. La validación
cualitativa de las etiquetas de los temas por dos analistas políticos confirmó que
25 de los 28 temas detectados (89.3\%) corresponden a temas electoralmente relevantes
y reconocibles en el contexto de la campaña peruana.

\textbf{Actividad 3: Evaluar sistema GraphRAG}
\\
Se diseñó un conjunto de 20 preguntas de evaluación sobre el grafo electoral, con
respuesta esperada verificada manualmente mediante consultas Cypher escritas por el
investigador. El sistema GraphRAG generó automáticamente las consultas Cypher
correspondientes a cada pregunta, obteniendo una exactitud del 85.0\% (17 de 20
consultas correctas). Los tres errores correspondieron a preguntas con comparaciones
temporales complejas entre más de dos períodos simultáneos.

\textbf{Actividad 4: Evaluar el sistema integrado}
\\
La evaluación final del sistema integrado comprendió la medición de la latencia
extremo a extremo del pipeline y la correlación de los indicadores ISN con encuestas
electorales disponibles. Se midió el tiempo total de procesamiento de 1,000 segmentos
del corpus en condiciones de carga media, obteniendo una latencia promedio de 2.3
segundos por segmento con la API de Gemini Flash. Esto permite al sistema procesar
en tiempo real los picos de actividad del ecosistema digital electoral (estimados en
500--800 publicaciones por minuto durante los debates presidenciales) con una
arquitectura de procesamiento paralelo de 4 workers simultáneos.

Para la validación de los indicadores ISN, se calculó la correlación de Pearson entre
los valores semanales del ISN por candidato generados por el sistema y los resultados
de tres encuestas electorales publicadas durante el mismo período, obteniendo un
coeficiente de $r = 0.74$ ($p < 0.05$), lo que indica una correlación fuerte y
estadísticamente significativa entre la percepción digital medida por el sistema y
la intención de voto declarada en las encuestas convencionales.

% ──────────────────────────────────────────────────────────────
\section{Despliegue}
% ──────────────────────────────────────────────────────────────

\textbf{Actividad 1: Diseñar dashboard de visualización}
\\
Se diseñó e implementó el dashboard interactivo del sistema de inteligencia electoral
utilizando el framework Streamlit de Python, desplegado en Google Cloud Run para
garantizar disponibilidad continua. El dashboard presenta tres vistas principales:
(1) la evolución temporal del ISN por candidato, (2) el mapa de calor de temas relevantes
por semana y (3) el módulo de consulta GraphRAG en lenguaje natural. La arquitectura
del prototipo del sistema integrado se muestra en la Figura \ref{5:fig24}.

% \begin{figure}[!ht]
%     \begin{center}
%         \includegraphics[width=0.95\textwidth]{4/figures/demo_flux.png}
%         \caption[Proceso de operación del prototipo del sistema de inteligencia electoral]{Proceso de operación del prototipo del sistema de inteligencia electoral: desde la ingesta de datos en tiempo real hasta la presentación de indicadores de percepción pública en el dashboard interactivo.\\
%             Fuente: Elaboración propia.}
%         \label{5:fig24}
%     \end{center}
% \end{figure}

\textbf{Actividad 2: Ejecutar prototipo con datos electorales en vivo}
\\
En esta actividad se ejecutaron las acciones de la arquitectura del prototipo
descritas en la figura anterior. La prueba piloto consistió en operar el sistema
durante una semana de campaña electoral activa, incluyendo un debate presidencial
que generó un pico de 3,200 tweets por minuto en el ecosistema digital. Durante esta
semana, el sistema procesó en tiempo real más de 48,000 publicaciones distribuidas
entre las cuatro fuentes de datos, generando los indicadores ISN, IRT e IVE para los
candidatos en liza. Los principales hallazgos de la prueba piloto se presentan en
la Figura \ref{5:fig25}.

% \begin{figure}[!ht]
%     \begin{center}
%         \includegraphics[width=0.92\textwidth]{4/figures/dashboard_resultados.png}
%         \caption[Resultados del dashboard de inteligencia electoral durante la semana de prueba piloto]{Resultados del dashboard de inteligencia electoral durante la semana de prueba piloto, mostrando la evolución del ISN por candidato, los temas emergentes identificados y la brecha entre cobertura mediática y sentimiento ciudadano.\\
%             Fuente: Elaboración propia.}
%         \label{5:fig25}
%     \end{center}
% \end{figure}

La prueba piloto confirmó la viabilidad del sistema para detectar en tiempo real
cambios bruscos en la percepción pública —como el pico de sentimiento negativo
registrado durante y después del debate presidencial— y para identificar los temas
que emergieron como consecuencia de los intercambios entre candidatos. El sistema
también identificó correctamente una discrepancia significativa entre la cobertura
positiva de un candidato en los medios digitales y el sentimiento negativo expresado
hacia ese mismo candidato en X/Twitter durante la misma semana, constituyendo un
hallazgo de inteligencia electoral relevante para los analistas usuarios del sistema.